% IEEE Transactions Journal Template
% For: IEEE Transactions on [Various Topics]
% Citation style: Numbered [1], [2], [3]
% Format: Two-column, Times New Roman 10pt
%
% Applicable journals:
%   - IEEE Transactions on Neural Networks and Learning Systems
%   - IEEE Transactions on Signal Processing
%   - IEEE Transactions on Image Processing
%   - IEEE Transactions on Pattern Analysis and Machine Intelligence
%   - IEEE Transactions on Industrial Electronics
%   - IEEE Transactions on Power Systems
%   - IEEE Transactions on Smart Grid
%   - IEEE Transactions on Robotics
%   - IEEE Transactions on Automation Science and Engineering
%   - IEEE Access (open access, no page limit)
%
% Usage:
%   pdflatex ieee_trans.tex
%   bibtex ieee_trans
%   pdflatex ieee_trans.tex
%   pdflatex ieee_trans.tex

\documentclass[journal]{IEEEtran}

% Essential packages
\usepackage[utf8]{inputenc}
\usepackage{amsmath,amssymb,amsfonts}
\usepackage{graphicx}
\usepackage{textcomp}
\usepackage{xcolor}
\usepackage{booktabs}
\usepackage{multirow}
\usepackage{cite}
\usepackage{algorithmic}
\usepackage{algorithm}
\usepackage{url}
\usepackage{hyperref}
\usepackage{siunitx}

% For double-column figures/tables
\usepackage{dblfloatfix}

% Correct bad hyphenation
\hyphenation{op-ti-cal net-works semi-conduc-tor}

\begin{document}

% ============================================================
% TITLE AND AUTHORS
% ============================================================
\title{Paper Title Here}

\author{%
  \IEEEauthorblockN{First Author\IEEEauthorrefmark{1},
  Second Author\IEEEauthorrefmark{1},
  Third Author\IEEEauthorrefmark{2}}\\
  \IEEEauthorblockA{\IEEEauthorrefmark{1}Department of Electrical Engineering,
    University Name, City, Country\\
    Email: \{first.author, second.author\}@university.edu}
  \IEEEauthorblockA{\IEEEauthorrefmark{2}School of Computer Science,
    University Name, City, Country\\
    Email: third.author@university.edu}
}

\markboth{IEEE Transactions on [Topic], Vol.~XX, No.~XX, Month~Year}%
{Author \MakeLowercase{\textit{et al.}}: Paper Title}

\maketitle

% ============================================================
% ABSTRACT (<=200 words, no references, no abbreviations)
% ============================================================
\begin{abstract}
Your abstract here. The abstract should be a concise summary of the
paper, typically 150--200 words. It should state the purpose of the
research, the methods used, the key results, and the main conclusions.
Do not cite references in the abstract. Define all abbreviations at
first use or avoid them entirely.
\end{abstract}

% ============================================================
% INDEX TERMS (IEEE controlled vocabulary)
% ============================================================
\begin{IEEEkeywords}
Keyword1, keyword2, keyword3, keyword4, keyword5.
\end{IEEEkeywords}

% ============================================================
% I. INTRODUCTION
% ============================================================
\section{Introduction}

The introduction should provide context, motivation, and a clear
statement of the paper's contributions. It should:

\begin{itemize}
  \item Motivate the problem and its significance
  \item Briefly survey related work (detailed review in Section~II)
  \item State the paper's main contributions clearly
  \item Outline the paper organization
\end{itemize}

% ============================================================
% II. RELATED WORK
% ============================================================
\section{Related Work}

Review the relevant literature, organized thematically. Compare and
contrast existing approaches, identify gaps, and position your work.

% ============================================================
% III. METHODOLOGY
% ============================================================
\section{Proposed Method}

\subsection{Problem Formulation}

Define the problem mathematically. Use IEEEeqnarray for aligned equations:
\begin{equation}
  \min_{x} f(x) \quad \text{subject to} \quad g(x) \leq 0
  \label{eq:objective}
\end{equation}

\subsection{Algorithm Description}

Describe the proposed algorithm. Use the algorithm environment:

\begin{algorithm}[t]
\caption{Algorithm Name}
\label{alg:algorithm}
\begin{algorithmic}[1]
  \REQUIRE Input data $X$, parameters $\theta$
  \ENSURE Output result $Y$
  \STATE Initialize $Y_0$
  \FOR{$t = 1$ to $T$}
    \STATE Compute $Y_t \leftarrow \text{Update}(Y_{t-1}, X, \theta)$
    \IF{convergence criterion met}
      \STATE \textbf{break}
    \ENDIF
  \ENDFOR
  \RETURN $Y_T$
\end{algorithmic}
\end{algorithm}

\subsection{Complexity Analysis}

Analyze the computational complexity:
\begin{itemize}
  \item Time complexity: $O(\cdot)$
  \item Space complexity: $O(\cdot)$
\end{itemize}

% ============================================================
% IV. EXPERIMENTAL SETUP
% ============================================================
\section{Experimental Setup}

\subsection{Datasets}

Describe datasets used, including size, source, and preprocessing.

\subsection{Baselines}

List baseline methods for comparison. Justify the selection.

\subsection{Implementation Details}

Report hardware, software, hyperparameters, and training details:
\begin{itemize}
  \item Hardware: GPU type, memory
  \item Software: Framework, version
  \item Hyperparameters: Learning rate, batch size, epochs
  \item Random seeds for reproducibility
\end{itemize}

\subsection{Evaluation Metrics}

Define metrics used (accuracy, precision, recall, F1, etc.).

% ============================================================
% V. RESULTS AND DISCUSSION
% ============================================================
\section{Results and Discussion}

\subsection{Main Results}

Present results in tables and figures. Use the IEEE table format:

\begin{table}[t]
  \centering
  \caption{Comparison with Baseline Methods}
  \label{tab:results}
  \begin{tabular}{lcc}
    \toprule
    Method & Metric 1 & Metric 2 \\
    \midrule
    Baseline 1 & 0.85 & 0.72 \\
    Baseline 2 & 0.88 & 0.75 \\
    \textbf{Ours} & \textbf{0.92} & \textbf{0.81} \\
    \bottomrule
  \end{tabular}
\end{table}

\subsection{Ablation Study}

Isolate the contribution of each component.

\subsection{Computational Cost}

Report training time, inference time, and memory usage.

% ============================================================
% VI. CONCLUSION
% ============================================================
\section{Conclusion}

Summarize the main contributions and findings. Discuss limitations
and outline directions for future work.

% ============================================================
% ACKNOWLEDGMENT
% ============================================================
\section*{Acknowledgment}

This work was supported by [Funding Agency] under Grant [Number].

% ============================================================
% REFERENCES (IEEE style - use IEEEtran bibliography style)
% ============================================================
\bibliographystyle{IEEEtran}
\bibliography{../references/references}

% ============================================================
% AUTHOR BIOGRAPHIES (required for IEEE Transactions)
% ============================================================
\begin{IEEEbiography}[{\includegraphics[width=1in,height=1.25in,clip,keepaspectratio]{author1.jpg}}]{First Author}
received the B.S. degree from University, City, Country, in Year,
and the Ph.D. degree from University, City, Country, in Year.
He/She is currently a [position] with the Department of [Field],
University, City, Country. His/Her research interests include
topic1, topic2, and topic3.
\end{IEEEbiography}

\begin{IEEEbiography}{Second Author}
received the B.S. degree from University, City, Country, in Year.
He/She is currently pursuing the Ph.D. degree with the Department
of [Field], University, City, Country. His/Her research interests
include topic1 and topic2.
\end{IEEEbiography}

\end{document}
